\section{Taxonomy}
\label{sec:taxonomy}

We organize the surveyed QIEA-for-MOO literature along five axes:
representation scheme, quantum-inspired operator type,
selection/diversity mechanism, hybridization with classical MOEAs, and
application domain. The axes were derived inductively from recurring
design choices across the literature identified in
Section~\ref{sec:introduction:scope}, rather than imposed from an
existing classification. They were not derived from the bibliometric
sweep of Section~\ref{sec:bibliometric}, but two of its keyword
clusters correspond to two of them: one gathers the operator and
convergence vocabulary of the operator axis, another the application
domains (Section~\ref{sec:bibliometric:clusters}).
Section~\ref{sec:critical-analysis} revisits which surveyed methods are
described in enough replicable detail to place along them with
confidence. The five axes are not independent -- a method's operator
type constrains which representations it can pair with, and its
selection mechanism is often inherited wholesale from the classical MOEA
it hybridizes with -- but each highlights a different design decision
worth comparing across methods, and Section~\ref{sec:feature-comparison}
uses them as the columns of its systematic comparison table. Before the
five axes, Section~\ref{sec:taxonomy:objectives} sets out the contrast
between single- and multi-objective QIEAs, which cuts across all of
them.

\subsection{Single- vs.\ Multi-Objective QIEAs}
\label{sec:taxonomy:objectives}

Single-objective QIEAs are not the subject of this survey, but they are
the baseline every QIEA-MOO method starts from, and comparing the two
shows which parts of a QIEA are tied to having one objective. Most of
the machinery carries over unchanged. The Q-bit chromosome, observation
into a binary string, and the rotation gate of Eq.~\ref{eq:qgate} are
the same in \citet{han2002qea}'s single-objective QEA and in
\citet{kim2006qmea}'s multi-objective QMEA, which applies them to a
multi-objective knapsack problem without changing the encoding.

What does not carry over is the rotation gate's \emph{reference
solution}. The gate rotates each Q-bit toward the corresponding bit of a
reference solution, and the single-objective QEA has an obvious choice:
the best solution found so far, kept per individual and refreshed by
local and global migration~\citep{han2002qea}. With several objectives
there is no single best solution, only a set of mutually non-dominated
ones, so every QIEA-MOO method must decide which of them each Q-bit
individual rotates toward. The QMEA/MQEA line of work answers by
updating each Q-bit individual toward non-dominated solutions in an
archive~\citep{kim2009mqea}, drawn at random, a step its authors call
global random migration~\citep{ryu2014dmqea}. This turns
the reference choice into a diversity mechanism in its own right:
different individuals are pulled toward different parts of the front
rather than all toward one point. \citet{li2007hqga}'s hybrid
quantum-inspired genetic algorithm for multi-objective flow-shop
scheduling gives a second answer: its Q-bit component scalarizes the
objectives with randomly drawn weights, so that for each weight vector
there is again a single best solution to rotate toward.
\citet{deng2024moqead}'s MOQEA/D gives a third: it decomposes the
problem into single-objective subproblems and assigns each Q-bit string
its own, so each string again has a single best solution. Outside the
thirteen methods, \citet{ho2013vectorqea} go further and redesign the
QEA for multiple objectives, with a fitness assignment that counts how
many objectives a solution improves and by how much, and a redesigned
rotation-angle update. For the
NSGA-II-hybridized methods
(Section~\ref{sec:taxonomy:hybridization}), the sources gathered for
this survey describe the outer NSGA-II loop in detail but do not, at
the level of detail we could confirm, state how the rotation gate's
reference solution is chosen.

Two smaller differences follow from the same cause. First, the
single-objective QEA keeps one reference solution per individual, while
a QIEA-MOO method must keep an archive of non-dominated solutions, which
adds memory and a per-generation archive-update cost. Second, the
single-objective literature treats Q-bit convergence as the goal:
\citet{han2004hepsilon} use the convergence of Q-bit probabilities
toward 0 or 1 as a termination criterion, and introduce the
H$_\epsilon$ gate specifically to keep those probabilities from
collapsing completely. In a multi-objective setting, a population
whose Q-bits all collapse toward one binary string has lost the
diversity it needs to cover a Pareto front, so a mechanism like the
H$_\epsilon$ gate's lower bound arguably matters more there than in the
single-objective case it was designed for. Section~\ref{sec:critical-analysis:so-to-mo}
returns to what these differences imply for the surveyed methods.

\subsection{Representation Scheme}
\label{sec:taxonomy:representation}

The representation axis asks what an individual encodes and how that
encoding is turned into an evaluable candidate solution. The canonical
QIEA representation is Han and Kim's Q-bit chromosome
(Section~\ref{sec:background:quantum}, Eq.~\ref{eq:qubit}): a string of
probability-amplitude pairs that is \emph{observed} into an ordinary
binary string before fitness evaluation, while the amplitudes themselves
carry the evolutionary search's memory across
generations~\citep{han2002qea}. This representation transfers directly
to combinatorial 0/1 problems -- \citet{kim2006qmea}'s QMEA observes a
Q-bit chromosome into the binary decision vector of a multi-objective
knapsack problem essentially unchanged from the single-objective case --
but requires adaptation for continuous decision spaces, where the
decision vector is real-valued rather than binary:
\citet{olvera2023continuous} keep observation-as-measurement unchanged
but add an explicit decoding step after it, converting each observed
binary string into floating-point decision variables before fitness
evaluation -- so that the precision with which a real value can be
represented is bounded by how many Q-bits are allotted to it. Feature subset selection sits at the
combinatorial end of this spectrum and is, per
\citet{vivek2025fss}'s survey, one of the domains where the Q-bit
representation is used closest to its original binary-observation form,
since a feature is naturally either included or excluded.

Notably, this survey's own case study (Section~\ref{sec:case-study}) does
not use a Q-bit representation at all: its individuals are classical,
domain-specific chromosomes -- lists of gate tuples
\texttt{(gate\_name, target\_qubit, ctrl\_qubit, angle)} -- evolved with
ordinary (non-quantum-inspired) genetic operators. This is not an
oversight but a reminder that the representation axis is orthogonal to
whether a method's \emph{problem domain} happens to be quantum computing;
we return to this distinction under application domain
(Section~\ref{sec:taxonomy:application}).

\subsection{Quantum-Inspired Operator Type}
\label{sec:taxonomy:operator}

The operator axis asks which quantum-mechanical phenomenon a method
draws its variation operator from. The dominant choice by far, across
essentially every method surveyed, is the rotation-gate operator of
Eq.~\ref{eq:qgate}, which nudges each Q-bit's amplitude pair toward
whichever basis state a lookup table favors given the current
individual's fitness relative to the best solution found so
far~\citep{han2002qea,vivek2025fss}. Beyond the rotation gate, two
further quantum phenomena recur, layered on top of rather than replacing
it: superposition, exploited by \citet{guzel2022qnsga2}'s QNSGA-II to
diversify sampling around the current population and reportedly speed
convergence relative to a plain NSGA-II baseline, and entanglement,
exploited by \citet{kumar2026eahqnsga2}'s EAH-QNSGA-II as an
entanglement-driven crossover operator that couples pairs of individuals
more tightly than independent single-point crossover would, combined
with an adaptive (rather than fixed-table) quantum rotation step and a
local-search refinement pass. The word
``superposition'' alone does not mark a method as quantum-inspired,
however: \citet{florea2023superposition}'s ``superposition of
populations'' runs NSGA-II and SPEA2 concurrently, each producing a
share of every generation set by its recent solution quality, with no
quantum-inspired representation or operator involved. Interference -- the third quantum
phenomenon named alongside superposition and entanglement in
Section~\ref{sec:introduction} as a general source of QIEA design
inspiration -- has no multi-objective method in the currently surveyed
set built explicitly around it as a variation operator. It does appear
on the single-objective side, and earlier than the rotation gate itself:
\citet{narayanan1996qiga}'s quantum-inspired genetic algorithm for the
travelling salesman problem, the earliest QIEA we cover, keeps several
parallel ``universes'' of solutions and builds new individuals with an
\emph{interference crossover} that takes successive genes from successive
universes. Interference was therefore part of quantum-inspired
evolutionary design from the start, but did not survive into the
rotation-gate-based multi-objective methods surveyed here.
Section~\ref{sec:critical-analysis:per-category} considers what this
narrow operator range implies for the field.

\subsection{Selection and Diversity Mechanism}
\label{sec:taxonomy:selection}

The selection axis asks how a method ranks and prunes its population
across objectives -- the same paradigms introduced as classical MOEA
baselines in Section~\ref{sec:background:moea}: crowding-distance-based,
reference-point-based, strength/density-based, decomposition-based, and
hypervolume-driven. Across the thirteen surveyed methods, selection is
more varied than the recent literature alone suggests. The three recent
NSGA-II-based methods inherit NSGA-II's selection \emph{unchanged}:
\citet{guzel2022qnsga2}'s QNSGA-II and \citet{kumar2026eahqnsga2}'s
EAH-QNSGA-II both retain its crowding distance verbatim, as do the three
continuous-space algorithms of \citet{olvera2023continuous}, applying
their quantum-inspired operators only to how new candidates are
generated, not to how the population is ranked or pruned. The older
QMEA/MQEA line ranks by Pareto dominance and keeps an external
non-dominated archive~\citep{kim2006qmea,kim2009mqea}, and MQEA-PS
fills that archive by user preference instead of
dominance~\citep{kim2011mqeaps}. \citet{zhang2012moqeavrp} maintain
diversity with a self-adaptive grid, in the spirit of PAES and
PESA~\citep{knowles2000paes,corne2000pesa}. \citet{deng2024moqead}'s
MOQEA/D decomposes the problem into single-objective subproblems, the
MOEA/D paradigm. \citet{li2007hqga} combine a randomly weighted sum of
the objectives, in their Q-bit component, with non-dominated sorting and
population trimming in a permutation-based component. The remaining
four methods do not state their selection mechanism in their
abstracts~\citep{wang2014powertrain,mousavi2019uav,deng2020gate,hussain2022moqiga}.
What none of the thirteen does is pair a Q-bit representation with
reference-point-based (NSGA-III-style), strength/density-based
(SPEA2-style), or hypervolume-driven (SMS-EMOA-style) selection, even
though the first two have recently gained the strongest theoretical
guarantees of their lineages
(Section~\ref{sec:critical-analysis:per-category}).
Section~\ref{sec:critical-analysis} treats this absence as a candidate
limitation rather than evidence that these paradigms are unsuitable.
This project's own case study pairs all three of NSGA-II, NSGA-III, and
SMS-EMOA selection with a classical (non-quantum-inspired) chromosome,
so it demonstrates the reverse pairing is viable, but on the domain-
specific representation side of the representation axis
(Section~\ref{sec:taxonomy:representation}), not the Q-bit side --
it does not by itself close this gap for QIEA-MOO specifically.

\subsection{Hybridization with Classical MOEAs}
\label{sec:taxonomy:hybridization}

Closely related to the selection axis, the hybridization axis asks how
tightly a method's quantum-inspired representation and operators are
grafted onto an existing classical MOEA's overall algorithmic skeleton,
versus how much of the surrounding framework is bespoke. On this axis,
the surveyed methods fall into three design points. In the first,
which the three most recent methods share, a method names and frames
itself as an explicit variant of a single classical MOEA -- NSGA-II in
every case found -- with the Q-bit
representation and one or two
quantum-inspired operators substituted into the generation step, while
the outer non-dominated-sorting-plus-crowding-distance loop is left
otherwise intact (\citet{guzel2022qnsga2},
\citet{kumar2026eahqnsga2}). \citet{olvera2023continuous} follow the
same pattern twice over: their proposed SQMEA is explicitly a
quantum-inspired version of NSGA-II, and NSGA-II (in both binary- and
floating-point-coded form) is also one of the classical baselines they
compare against.
A second design point, older and more common among the most-cited
methods, extends a QIEA with its own multi-objective machinery rather
than grafting it onto a named MOEA: \citet{kim2006qmea}'s QMEA and its
successors~\citep{kim2009mqea,kim2011mqeaps} build Pareto ranking and an
archive directly into the QEA, \citet{zhang2012moqeavrp} add an adaptive
grid, and \citet{deng2024moqead} borrow MOEA/D's decomposition idea
without presenting the result as a variant of MOEA/D itself. A third
design point hybridizes the QIEA with another metaheuristic rather than
with an MOEA:
\citet{li2007hqga} pair a Q-bit genetic algorithm with a
permutation-based genetic algorithm, the first exploring a binary space
and the second refining schedules directly, and \citet{deng2020gate}
combine a QEA with particle swarm optimization and a niche co-evolution
strategy. These two are also the most-cited multi-objective QIEAs in
the bibliometric corpus, so the NSGA-II-variant design point, although
the most common among recent methods, is also the newest of the three
and so far the least cited. This survey's own case study
(Section~\ref{sec:case-study}) sits at the opposite end of the same
axis from the QIEA-MOO literature above: it hybridizes classical MOEAs
(NSGA-II, NSGA-III, SMS-EMOA, used as unmodified per-block optimizers)
with a domain-specific, non-quantum-inspired chromosome, rather than
hybridizing a classical MOEA with a Q-bit representation -- the same
"reverse pairing" noted under selection above, presented here as a
structural contrast rather than as an example of QIEA-MOO hybridization
itself.

\subsection{Application Domain}
\label{sec:taxonomy:application}

The application-domain axis asks what kind of problem a method is
evaluated on. Two largely disjoint clusters emerge across the literature
surveyed here. The first cluster comprises classical combinatorial or
engineering optimization problems to which a quantum-\emph{inspired}
algorithm is applied on ordinary classical hardware: multi-objective 0/1
knapsack~\citep{kim2006qmea,lu2013apmqea}, continuous and DTLZ benchmark
problems~\citep{olvera2023continuous,kim2011mqeaps}, flow-shop and
cloud-workflow scheduling~\citep{li2007hqga,hussain2022moqiga}, vehicle
routing~\citep{zhang2012moqeavrp}, airport gate
allocation~\citep{deng2020gate,deng2024moqead}, robot path
planning~\citep{kim2011mqeaps}, UAV task
allocation~\citep{mousavi2019uav}, electric-vehicle powertrain
design~\citep{wang2014powertrain}, and electric-vehicle
charging-station siting~\citep{kumar2026eahqnsga2}. In every one
of these, ``quantum'' describes the algorithm's inspiration, not the
problem: knapsacks, schedules, routes, airport gates, vehicles, and
charging networks are entirely classical objects. The second cluster is the
opposite pairing: quantum circuit optimization, where the problem
domain is genuinely quantum (a target unitary or circuit to reproduce)
but the optimizer applied to it, in this survey's own case study and its
foundational reference~\citep{ghlib2026scalable}, is an unmodified
classical MOEA rather than a quantum-inspired one. The two do meet on
the single-objective side: the bibliometric corpus of
Section~\ref{sec:bibliometric} contains QIEAs that evolve quantum
circuits at the gate level~\citep{ding2008qcircuits}, design or
synthesize circuits for given gate
libraries~\citep{krylov2019qeqea,chou2022qcsynthesis}, and search
parameterized-circuit architectures for quantum neural
networks~\citep{li2025aqeaqas}. We did not, however, find a
\emph{multi-objective} QIEA applied to a quantum-computing problem
domain, even though such problems are naturally multi-objective --
fidelity against depth and gate count in this survey's own case study.
Whether that combination would offer any advantage over either cluster
alone is an open question we return to as a candidate future direction
in Section~\ref{sec:open-challenges}.
